\documentclass[12pt,a4paper]{ctexrep}

% --- 宏包 ---
\usepackage{geometry,setspace,titlesec,fancyhdr}
\usepackage{amsmath,booktabs,tabularx,graphicx}
\usepackage{enumitem}
\usepackage[hidelinks]{hyperref}
\usepackage{tikz}
\numberwithin{equation}{section}
\usetikzlibrary{shapes.geometric, arrows, positioning, calc}

% --- 页面边距与布局 ---
\geometry{top=2.5cm,bottom=2.54cm,left=2.5cm,right=2.5cm}

% --- 字体设置 ---
\setCJKmainfont{SimSun}
\setCJKsansfont{SimHei}
\setmainfont{Times New Roman}
\newcommand{\song}{\songti}
\newcommand{\hei}{\heiti}

% --- 段落与页眉页脚 ---
\setlength{\parskip}{0pt}
% --- 页眉页脚设置 ---
\pagestyle{fancy}
\fancyhf{}
\cfoot{\song\footnotesize\thepage}
\renewcommand{\headrulewidth}{0pt}


% --- 重新定制标题格式 ---
\titleformat{\section}[block]{\hei\zihao{-3}}{\thesection}{1em}{}
\titlespacing{\section}{0pt}{8pt}{8pt}

\titleformat{\subsection}[block]{\hei\zihao{4}}{\thesubsection}{1em}{}
\titlespacing{\subsection}{0pt}{5pt}{5pt}

\titleformat{\subsubsection}[block]{\hei\zihao{4}}{\thesubsubsection}{1em}{}
\titlespacing{\subsubsection}{0pt}{5pt}{5pt}

\renewcommand{\thesection}{\arabic{section}}
\renewcommand{\thesubsection}{\thesection.\arabic{subsection}}
\renewcommand{\thesubsubsection}{\thesubsection.\arabic{subsubsection}}
\setcounter{secnumdepth}{3}

% --- 引用符号定制（上标、带中括号） ---
\newcommand{\citeup}[1]{\textsuperscript{[#1]}}


\begin{document}
	
	% ===================================================
	% 前置部分：中文摘要 ===================================================
	\newpage
	\pagenumbering{arabic} % 设置阿拉伯数字页码
	
	% --- 1. 中文摘要 ---
	\begin{center}
		{\hei\zihao{2}标题}
	\end{center}
	\noindent{\hei\zihao{-3}摘要}\par
	\vspace{11pt}
	
	\begin{spacing}{1.0}
		{\song\zihao{4} \setlength{\parindent}{2em}
			中文摘要
		}
	\end{spacing}
	
	\vspace{1em}
	\noindent{\hei\zihao{4}关键词：1；2；3；4}
	
	
	%===================================================
	% 正文部分
	%===================================================
	\begin{spacing}{1.10}
		{\song\zihao{-4}
		%引言部分
			\section{引言}
			%简要概括这个研究目标的重要性
			
			%问题重述：
			%简要概括问题的具体是什么，转化数学建模目标
			
			\section{问题分析}
			%承接上文的建模目标，分析这种情况不应该只考虑那种模型，还应引入什么
			%在此处插入主流程图
			
			\section{问题假设}
			%问题的重要假设，支持后续建模
			
			\section{符号即标准说明}
			%文章中的重要符号
			
		%建模准备工作
			\section{数据预处理}
			\subsection{数据总览}
			%简要说一下数据集都有什么项目
			\subsection{数据理解与异常值处理}
			%简要说一下一些不太好理解的数据应该怎么定义，如何参与建模
			%在此基础上，对异常值应该怎么处理
			\subsection{数据统一与变换}
			%将数据统一或者变换，便于后续建模
			\subsection{训练集与测试集设计}
			%讲解一下如何防止数据泄露，防止污染
			\subsection{建模准备总结}
			%总结数据预处理，衔接下文正式建模
			%在此插入建模准备流程图
		
		%正式建模	
			\section{问题1建模}
			\subsection{预实验与文献综述}
			%对数据中进行简要分析，选择出主模型，结合文献给出主模型的合理性
			\subsection{主模型建立}
			%用数学公式表达模型
			%指出这个模型解决了什么问题
			\subsection{主模型求解}
			%给出模型结果，说明求解过程采用什么算法
			%求解的参数如何进行优化
			\subsection{误差分析与稳健性检验}
			%说明误差结构来源
			%对参数进行稳健性分析
			\subsection{第一问结论}
			%总结上文，给出那些及结论可以保留并参与后续建模。你哪些结论因为什么原因应当慎重考虑
			%注意正面回应题目的要求，如果题目要求形成可以参考的表格或者分界点应明确用表格给出（短表直接写，长表放到appendix_data目录）
			
			\section{结论}
			%文章主结论，最后正面回应题目要求
		
		%模型评价与推广	
			\section{模型推广与评价}
			\subsection{模型优点}
			%正确处理了什么问题，有什么创新行性
			\subsection{模型局限}
			%在什么地方上，由于什么原因，没有达到
			\subsection{模型推广}
			%可以参考什么，哪些应该慎重，
		}
	\end{spacing}
	
	\section{AI使用声明}
	
	% 参考文献部分
	\section{参考文献}	
	\begin{spacing}{1.25} % 1.25 倍行距
		\zihao{5}             % 五号字
	\begin{enumerate}[label={[\arabic*]}, leftmargin=2em, labelwidth=1.5em, labelsep=0.5em, itemsep=0pt, parsep=0pt, partopsep=0pt, topsep=0pt]		
		\item 						
	\end{enumerate}
	\end{spacing}
	
	%附录部分
	%创建appendix文档 
	
	%附录A 目录
	%创建appendix_catelog 附录部分目录
	
	%附录B 直接引用数据
	%appendix_data_citation
	%论文直接引用的重要数据，比如预实验设计和结果
	
	%附录C 重要长表数据
	%appendix_data_adition
	%部分划分结果，比如数据清洗报告等重要但是没有直接引用的数据
	
	%附录D 支撑材料目录
	%appendix_contents 支撑文件目录，给出文件名和用途
	
	%附录E  代码文件
	%appendix_code     代码文件
\end{document}

